\section{Results}

Every figure in this section is generated by the code in \texttt{src/}, and every
number quoted is produced by running it. No figure is drawn by hand. The parameter
values are illustrative rather than fitted: their role is to establish the qualitative
structure and the orders of magnitude, not to describe any particular patient.

%---------------------------------------------------------------
\subsection{Model A: the degeneracy and its resolution}

\begin{figure}[htbp]
\centering
\includegraphics[width=\textwidth]{fig_A1_channels.pdf}
\caption{Left: the two elimination channels. The enzymatic term saturates at
$\theta$ while the alternative term does not, and this difference in functional
form is the only thing distinguishing them. In the shaded band, where clinically
stable patients are maintained, both are approximately linear. Right: two parameter
sets that produce identical total clearance at the operating point and diverge
only well outside it. These two patients differ threefold in enzymatic capacity and
are indistinguishable from measurements taken where they live.}
\label{fig:A1}
\end{figure}

\begin{figure}[htbp]
\centering
\begin{minipage}[t]{0.48\textwidth}
\centering
\includegraphics[width=\textwidth]{fig_A2_conditioning.pdf}
\end{minipage}\hfill
\begin{minipage}[t]{0.48\textwidth}
\centering
\includegraphics[width=\textwidth]{fig_A5_load_error.pdf}
\end{minipage}
\caption{Left: condition number of the column-normalised sensitivity matrix against
the relative excursion of the metabolite, on log axes. The slope is $-2$: conditioning
degrades quadratically as the patient is stabilised, and a patient varying by a few
per cent between visits presents a problem conditioned at $10^{5}$ to $10^{6}$.
Right: an error in the assumed substrate load propagates into $\theta$ with an
amplification factor of approximately $1.4$, because the enzymatic term is recovered
as a difference of two larger quantities.}
\label{fig:A2}
\end{figure}

\begin{figure}[htbp]
\centering
\includegraphics[width=\textwidth]{fig_A3_indistinguishable.pdf}
\caption{Left: a selective perturbation renders the two patients of Figure~\ref{fig:A1}
formally distinguishable. The perturbation shown is carnitine, which acts on the
alternative channel $k_e$ and is administered to every patient in the accessible cohort.
The separation remains near $1\%$ across the plausible dose range, well below assay
variability: formal identifiability and practical identifiability part company here.
Right: the separation plotted against the excursion of $m$ that the perturbation actually
achieves, for both channels. The two curves lie on top of one another, which is
Proposition~\ref{prop:channel} made visible. The information is bought by moving the
patient, and it does not matter which channel is used to move them. This is what allows
the analysis to be conducted on the input that exists rather than the one that does not.}
\label{fig:A3}
\end{figure}

\begin{figure}[htbp]
\centering
\includegraphics[width=\textwidth]{fig_A4_population.pdf}
\caption{Left: at steady state a cohort of $N$ patients sharing the population
parameters yields $N$ equations in $N+2$ unknowns. The deficiency is two for every
$N$: enlarging the cohort adds an equation and an unknown simultaneously and never
helps. Right: with two anchors the system closes exactly. A null genotype supplies
$k_e$ directly, since its equation loses the enzymatic term entirely; one patient
observed at two distinct loads then supplies $K_m$; every remaining $\theta_i$
follows from a single equation. Recovery is exact in the noise-free case.}
\label{fig:A4}
\end{figure}

\begin{proposition}[Anchor condition]
Let $N$ patients at steady state share $(K_m, k_e)$. If the cohort contains at least
one patient with a null genotype, so that $\theta_j = 0$, and at least one patient
observed at two distinct steady states, then $K_m$, $k_e$ and every $\theta_i$ are
uniquely determined.
\end{proposition}

Both anchors are present in the Saudi series: the frameshift variant \emph{PCCB}
c.990dupT occurs in seven families and is a null allele~\cite{saudi84}, and documented changes in % C26
prescribed protein intake supply the second condition.

Two things justify setting $\theta_j = 0$ for the null patients, and neither is a claim
about residual carboxylase activity. The first is molecular: c.990dupT is a frameshift
introducing a premature termination codon, and transcripts carrying one are subject to
nonsense-mediated decay, so no functional PCC is translated. % C32
The second is structural, and it is the stronger of the two because it does not depend on
the first. The construction does not require that \emph{no} propionate leaves by any other
route. It requires only that whatever else eliminates propionate is not
$\theta\, g(m)$. Any residual elimination, of whatever molecular origin, enters
\eqref{eq:A} as part of the alternative channel and is absorbed into $k_e$, which is
precisely the quantity the null patients measure. % C27
An earlier version of this note argued instead that residual carboxylation in PCC-deficient
cells derives from a distinct carboxylase. That may well be true, but no source for it was
found, and the anchor does not need it.

\clearpage
%---------------------------------------------------------------
\subsection{Model B: the load ceiling and what it costs}

\begin{figure}[htbp]
\centering
\begin{minipage}[t]{0.49\textwidth}
\centering
\includegraphics[width=\textwidth]{fig_B1_ceiling.pdf}
\end{minipage}\hfill
\begin{minipage}[t]{0.49\textwidth}
\centering
\includegraphics[width=\textwidth]{fig_B2_bifurcation.pdf}
\end{minipage}
\caption{Left: with both channels saturating, total clearance is bounded by
$\theta + k_e$. Loads below the ceiling admit exactly one steady state; loads above it
admit none. Right: the equilibrium as a function of load. It escapes to infinity as the
load approaches the ceiling rather than colliding with a second branch, so this is not
a saddle-node. The clinically operative threshold arrives earlier, when the equilibrium
migrates past the tolerable level.}
\label{fig:B1}
\end{figure}

\begin{figure}[htbp]
\centering
\begin{minipage}[t]{0.47\textwidth}
\centering
\includegraphics[width=\textwidth]{fig_B3_thetacrit.pdf}
\end{minipage}\hfill
\begin{minipage}[t]{0.50\textwidth}
\centering
\includegraphics[width=\textwidth]{fig_B4_cost.pdf}
\end{minipage}
\caption{Left: $\thc$ over the plane of the two quantities it depends on, neither of
which is measured. The critical capacity is an operational definition rather than a
physiological constant, and the figure shows how much of its value is inherited from
that choice. Right: the price of biological fidelity. Once both channels saturate they
share a functional form and are separated only by differing half-saturation constants.
At a ten per cent excursion the condition number rises from $1.1 \times 10^{5}$ in
Model A to $4.8 \times 10^{7}$ in Model B, a factor of roughly $450$.}
\label{fig:B2}
\end{figure}

\begin{remark}
This is the central trade-off of the modelling exercise and should be stated plainly:
\textbf{the more faithful model is the harder one to identify}. Model A is retained as
the working model for the identifiability analysis not because it is more accurate,
which it is not, but because it represents the most favourable case. If identification
fails in Model A it fails a fortiori in Model B.
\end{remark}

\clearpage
%---------------------------------------------------------------
\subsection{Model C: depletion, a testable prediction, and a well-posed threshold}

A two-state depletion model with constant replenishment does not in fact produce
bistability. Setting $\dot c = 0$ in $\dot c = R - k c m$ gives a quasi-steady flux of
exactly $R$, which is saturation and not decline. Decline requires that replenishment
itself be impaired as the metabolite accumulates. The mechanism is already present in
the biology: sequestration of free coenzyme A causes mitochondrial dysfunction, which
degrades the cell's capacity to sustain carnitine handling. We therefore take
$R(m) = R_0/(1 + m/M_i)$, and the resulting effective clearance rises, peaks and falls.

\begin{figure}[htbp]
\centering
\begin{minipage}[t]{0.49\textwidth}
\centering
\includegraphics[width=\textwidth]{fig_C1_phaseline.pdf}
\end{minipage}\hfill
\begin{minipage}[t]{0.49\textwidth}
\centering
\includegraphics[width=\textwidth]{fig_C2_fold.pdf}
\end{minipage}
\caption{Left: the phase line under Model C. At baseline load the system has two
equilibria, a stable attractor at low metabolite and an unstable one at high
metabolite which acts as a separatrix. Under peak load neither exists. Right: the
corresponding bifurcation diagram. The two branches approach and collide, which is a
genuine saddle-node and distinguishes Model C from Model B.}
\label{fig:C1}
\end{figure}

\begin{figure}[htbp]
\centering
\begin{minipage}[t]{0.48\textwidth}
\centering
\includegraphics[width=\textwidth]{fig_C3_timeseries.pdf}
\end{minipage}\hfill
\begin{minipage}[t]{0.50\textwidth}
\centering
\includegraphics[width=\textwidth]{fig_C4_basin.pdf}
\end{minipage}
\caption{Left: the discriminating prediction. An identical transient illness is applied
to both models. Under Model B the metabolite returns to baseline once the illness
resolves. Under Model C the patient has crossed the separatrix and continues to
deteriorate at baseline load. Right: the basin of attraction. There is no upper stable
state, so the clinically meaningful object is the separatrix rather than a hysteresis
loop. Reducing the load raises the separatrix and can recapture the patient, which is
why suppression of catabolism by intravenous glucose is the mainstay of treatment;
lowering the metabolite directly, by haemofiltration, is the alternative route when the
load cannot be reduced far enough.}
\label{fig:C2}
\end{figure}

\begin{remark}[The discriminant, and what it cannot be measured in]
Models B and C differ in a prediction: under B, returning $s$ to baseline returns $m$ to
baseline, and under C it does not, because the separatrix at baseline load still lies below
the patient. That difference is a statement about the model's state $m$, and the obstacle
is that $m$ cannot be watched through a crisis. Deep phenotyping of acute decompensation
found propionylcarnitine, propionylglycine and propionic acid essentially unaltered during
the event, while the C2/C3 ratio fell about twenty-fold~\cite{amd_phenotyping}. % C24
The single marker does not track the state through the excursion the discriminant is about,
so ``does 3-HP return to baseline'' is not the question to take to the records. Two
quantities that do move are: the acetylcarnitine to propionylcarnitine ratio, which falls
sharply and recovers, and the within-patient fall in free carnitine relative to the same
patient's baseline, which is what the depletion mechanism of Model C actually asserts.
\end{remark}

\begin{figure}[htbp]
\centering
\includegraphics[width=\textwidth]{fig_C5_margin.pdf}
\caption{The intrinsic critical capacity, Proposition~\ref{prop:intrinsic}, computed at
baseline load. Left: the attractor and the separatrix as functions of the residual
capacity $\theta$, on a logarithmic metabolite axis. As $\theta$ falls the attractor rises
and the separatrix falls, and at $\thc^{\,\mathrm{int}} = 167$ they merge in a saddle-node;
below that value no equilibrium exists at all and the patient cannot be maintained even at
baseline load. Nothing in this construction refers to a tolerable metabolite level: the
threshold is a property of the model's equilibrium structure, not of a clinical convention.
Right: the distance between the two, which is the safety margin. It is the quantity the
clinical decision actually turns on, since it says how far the patient can be pushed before
the crisis becomes self-sustaining, and it closes continuously to zero at the threshold. At
the nominal $\theta = 300$ the margin is $398$. Neither Model A nor Model B admits this
figure: A never loses its equilibrium, and B loses one to infinity rather than to a
collision, so in neither is there a second equilibrium to measure a margin against.}
\label{fig:C5}
\end{figure}

\begin{remark}[Why this is the contribution of Model C]
Model C is usually introduced as the most biologically faithful of the three, and that
framing undersells it. Its distinguishing property is not fidelity but well-posedness: it
is the only one of the three models in which $\thc$ can be defined without importing a
convention. In A and B the threshold is an operational choice, and
Figure~\ref{fig:B2}~(left) shows how much of its numerical value is inherited from that
choice. In C it is the location of a saddle-node.
\end{remark}

\begin{remark}[A second, independent check]
The depletion mechanism predicts that free carnitine should fall sharply during
decompensation relative to the same patient's baseline. Free carnitine and the
acylcarnitine to free-carnitine ratio are measured routinely in these patients, so this
too requires no new experiment.
\end{remark}
